\documentclass[a4paper]{article}
\usepackage[italian]{babel}
\usepackage[margin=16mm]{geometry}
\usepackage{tikz}
\usetikzlibrary{shapes.geometric}

\tikzset{
  entita/.style={
    rectangle, draw, semithick, minimum width=2.65cm,
    minimum height=8mm, font=\bfseries
  },
  relazione/.style={
    diamond, draw, semithick, aspect=2.4, inner sep=1.5pt,
    minimum width=2.2cm, minimum height=9mm, font=\small
  },
  attributo/.style={inner xsep=3pt, inner ysep=2pt, font=\small},
  linea attributo/.style={draw=black!65, thin},
  cardinalita/.style={font=\small, inner sep=2pt}
}

\begin{document}

\begin{center}
  \Large\bfseries Diagramma entità--relazioni di Fixd
\end{center}

Il modello rappresenta gli account, gli annunci e le categorie, oltre ai
calendari dei professionisti e alle relative fasce di disponibilità.
Ogni annuncio appartiene a una categoria e rimanda al calendario del
suo autore attraverso l'account.

\vspace{8mm}
\begin{center}
\begin{tikzpicture}
  % Entità e relazioni
  \node[entita] (utente) at (0,0) {Utente registrato};
  \node[relazione] (pubblica) at (3.65,0) {Pubblica};
  \node[entita] (annuncio) at (7.3,0) {Annuncio};

  \node[relazione] (possiede) at (0,-1.85) {Possiede};
  \node[entita] (calendario) at (0,-3.7) {Calendario};
  \node[relazione] (include) at (0,-5.55) {Include};
  \node[entita] (disponibilita) at (0,-7.4) {Disponibilità};

  \node[relazione] (classifica) at (7.3,-1.85) {Classifica};
  \node[entita] (categoria) at (7.3,-3.7) {Categoria};

  % Cardinalità minime e massime dal lato dell'entità.
  \draw (utente) -- node[cardinalita, above, midway, yshift=4.5mm] {$(0,N)$} (pubblica);
  \draw (annuncio) -- node[cardinalita, above, midway, yshift=4.5mm] {$(1,1)$} (pubblica);
  \draw (utente) -- node[cardinalita, left, pos=.28] {$(1,1)$} (possiede);
  \draw (calendario) -- node[cardinalita, left, pos=.28] {$(1,1)$} (possiede);
  \draw (calendario) -- node[cardinalita, left, pos=.28] {$(0,N)$} (include);
  \draw (disponibilita) -- node[cardinalita, left, pos=.28] {$(1,1)$} (include);
  \draw (annuncio) -- node[cardinalita, right, pos=.28] {$(1,1)$} (classifica);
  \draw (categoria) -- node[cardinalita, right, pos=.28] {$(0,N)$} (classifica);

  % Ogni riga collega un solo attributo al bordo della sua entità.
  \foreach \etichetta/\quota/\attacco in {
    {\underline{ID utente}}/1.9/3mm,
    Nome/1.14/1.8mm,
    Cognome/.38/.6mm,
    Email/-.38/-.6mm,
    {Hash password}/-1.14/-1.8mm,
    Telefono/-1.9/-3mm
  } {
    \draw[linea attributo] ([yshift=\attacco]utente.west)
      -- (-2.55,\quota) node[attributo, anchor=east] {\etichetta};
  }
  \foreach \etichetta/\quota/\attacco in {
    {\underline{ID annuncio}}/1.9/3mm,
    Titolo/1.14/1.8mm,
    Descrizione/.38/.6mm,
    Tariffa/-.38/-.6mm,
    Latitudine/-1.14/-1.8mm,
    Longitudine/-1.9/-3mm
  } {
    \draw[linea attributo] ([yshift=\attacco]annuncio.east)
      -- (9.85,\quota) node[attributo, anchor=west] {\etichetta};
  }

  \draw[linea attributo] (calendario.east)
    -- (2.55,-3.7) node[attributo, anchor=west] {\underline{ID calendario}};
  \draw[linea attributo] ([yshift=1.5mm]categoria.east)
    -- (9.85,-3.32) node[attributo, anchor=west] {\underline{ID categoria}};
  \draw[linea attributo] ([yshift=-1.5mm]categoria.east)
    -- (9.85,-4.08) node[attributo, anchor=west] {Nome};

  \foreach \etichetta/\quota/\attacco in {
    {\underline{ID fascia}}/-6.26/2.7mm,
    Giorno/-7.02/.9mm,
    {Ora inizio}/-7.78/-.9mm,
    {Ora fine}/-8.54/-2.7mm
  } {
    \draw[linea attributo] ([yshift=\attacco]disponibilita.east)
      -- (2.55,\quota) node[attributo, anchor=west] {\etichetta};
  }
\end{tikzpicture}
\end{center}
\vspace{8mm}

\small
\noindent\textbf{Legenda.} Gli attributi sottolineati identificano le entità.
La coppia $(\min,\max)$ indica quante volte un'istanza dell'entità adiacente
può partecipare alla relazione. $N$ indica un massimo di più istanze.

\smallskip
\noindent La ricerca è pubblica. Per pubblicare o contattare un professionista
occorre accedere al proprio account. Il modello comprende i dati persistenti
descritti nel progetto. La ricerca, i suggerimenti AI e il contatto sono
funzionalità dell'applicazione.

\end{document}
